\scanpage{148}
% archive foliation: 147

cette catégorie est semi-simple). Si donc \(\mathcal C\) est
une catégorie abélienne, les foncteurs additifs
de \(\mathcal M\) dans \(\mathcal C\) sont autom. exacts. Cela
prouve que le ``Pb universel'' \inserted{(avec un grain de sel)}
en termes
de catégories abéliennes (Variante \(2\)) est
résolu par \(\mathcal M\).

Prouvons donc que \(\mathcal M\) est une catégorie
\uline{abélienne} \uline{semi-simple}. Cela équivaut à ceci:

\deleted{tout morphisme est \unsure{dir}}
\unsure{(i)|(iv)} Tout objet \deleted{\(\unsure{X}\)} satisfait condition des chaînes ascendantes,
et descendantes pour les \unsure{ss-objets|sous-objets} facteurs directs,

\deleted{i.e. \unsure{l'on a|l'est} des \unsure{pr}}
(ii) Tout morphisme \(u:P\longrightarrow Q\) est ``direct'',
i.e. isom à un morphisme
\[
A\times B\xrightarrow[\operatorname{inj}_1\circ\operatorname{pr}_2]{}B\times C
\]

Pour vérifier (ii), \deleted{\unsure{considérons}} considérons \(R=P\times Q\),
et l'endomorphisme \(\deleted{\unsure{pr}}\ v=\operatorname{inj}_2u\operatorname{pr}_2\) de \(R\).
On voit aisément que pour \deleted{\unsure{que}} (ii) \inserted{soit vrai}
pour \(u\),
il \inserted{\(^{(*)}\)} suffit \deleted{de} qu'il le soit pour \(v\), donc on
est ramené au cas où \(P=Q\). D'autre part,
soit \deleted{\(\unsure{R}\)} \(S\) un autre objet, et sur
\(R=P+S\), \(v=u\oplus0_S\), alors on constate que la
propriété est vraie pour \(u\) ssi elle l'est pour \(v\). Cela nous

\margin{\((*)\) dans un \unsure{carré} commut.
il est innocent de modifier
\(P\longrightarrow Q\) par
\(P\unsure{\in|\subset}P',\ Q\unsure{\in|\subset}Q'\)
\diagram{
\begin{tikzcd}[column sep=small,row sep=small]
&&Q'\\
{} \arrow[r] & P \arrow[r] & Q \arrow[u]
\end{tikzcd}
}}
